% Archived pseudocode, removed from main.tex on the grounds that every step is
% already covered by prose + formal definitions in Section III (Parts B and E),
% so keeping both as boxed \begin{algorithm} environments was redundant in the
% short (conference) version. Kept here for a future long/journal version where
% an explicit algorithmic presentation may still be wanted alongside the
% definitions, e.g. as an implementation appendix.
%
% Requires (already loaded by main.tex): \usepackage[linesnumbered,ruled,vlined]{algorithm2e}
%
% To reuse: % Archived pseudocode, removed from main.tex on the grounds that every step is
% already covered by prose + formal definitions in Section III (Parts B and E),
% so keeping both as boxed \begin{algorithm} environments was redundant in the
% short (conference) version. Kept here for a future long/journal version where
% an explicit algorithmic presentation may still be wanted alongside the
% definitions, e.g. as an implementation appendix.
%
% Requires (already loaded by main.tex): \usepackage[linesnumbered,ruled,vlined]{algorithm2e}
%
% To reuse: % Archived pseudocode, removed from main.tex on the grounds that every step is
% already covered by prose + formal definitions in Section III (Parts B and E),
% so keeping both as boxed \begin{algorithm} environments was redundant in the
% short (conference) version. Kept here for a future long/journal version where
% an explicit algorithmic presentation may still be wanted alongside the
% definitions, e.g. as an implementation appendix.
%
% Requires (already loaded by main.tex): \usepackage[linesnumbered,ruled,vlined]{algorithm2e}
%
% To reuse: % Archived pseudocode, removed from main.tex on the grounds that every step is
% already covered by prose + formal definitions in Section III (Parts B and E),
% so keeping both as boxed \begin{algorithm} environments was redundant in the
% short (conference) version. Kept here for a future long/journal version where
% an explicit algorithmic presentation may still be wanted alongside the
% definitions, e.g. as an implementation appendix.
%
% Requires (already loaded by main.tex): \usepackage[linesnumbered,ruled,vlined]{algorithm2e}
%
% To reuse: \input{appendix-algorithms} from within a section, or copy the
% relevant \begin{algorithm}...\end{algorithm} block directly. Labels are kept
% as in the original (alg::acl2use, alg::checking); re-check for label clashes
% before including both this file and any other file using the same labels.

% ---------------------------------------------------------------------------
% Originally in Part B (sec:acl2use, "Translating ACL into a Class Diagram"),
% introduced by: "Algorithm~\ref{alg::acl2use} summarizes the translation and
% its compatibility rule: for each concrete role pair $(R_1, R_2)$ sharing a
% group scope, no invariant is generated if they can never share an instance
% or if $\mathit{Compat}$ declares them compatible over that scope; otherwise
% \texttt{inv NoConflict\_$R_1$\_$R_2$} is generated on \texttt{Agent}."
% ---------------------------------------------------------------------------
\begin{algorithm}
\small
\ttfamily
\caption{Translating an ACL model into a USE class model}
\label{alg::acl2use}
\KwIn{an ACL model $(E, R, G, \mathit{Assoc}, \mathit{Compat})$}
\KwOut{a USE \texttt{.use} class model}

create a USE class for every entity in $E$, role in $R$, and group in $G$\;
translate role extension into USE generalization\;
\ForEach{relation in $\mathit{Assoc}$}{
  add the corresponding USE association/aggregation/composition, with the same cardinalities\;
}
create class \texttt{Agent}\;
\ForEach{ownership relation between a group $g$ and a role $r$}{
  add association \texttt{Agent\_plays\_$r$} between \texttt{Agent} and $r$\;
  add association \texttt{$r$\_in\_$g$} between $g$ and $r$, with cardinality $\mathit{card}(g,r)$\;
}
\ForEach{ownership relation between a parent group and a child group}{
  add a composition between them\;
}
\ForEach{pair of concrete roles $(R_1,R_2)$ that can share a group scope}{
  \eIf{$(R_1,R_2) \in \mathit{Compat}$ over that scope}{
    generate no invariant\;
  }{
    generate \texttt{inv NoConflict\_$R_1$\_$R_2$} on \texttt{Agent}\;
  }
}
\Return the generated \texttt{.use} class model\;
\end{algorithm}

% ---------------------------------------------------------------------------
% Originally in Part E (sec:checking, "Checking i* and BPMN over State
% Transitions"), introduced by: "Algorithm~\ref{alg::checking} restates this
% as a procedure," following the four failure-list definitions
% (aclFailures, bpmnFailures, goalFailures, oracleFailures).
% ---------------------------------------------------------------------------
\begin{algorithm}
\small
\ttfamily
\caption{Checking i* against BPMN on a checkpoint trace}
\label{alg::checking}
\KwIn{a USE class model generated from ACL, an initial SOIL script compiled from AOL, an i* model, an ISCN oracle, a BPMN model}
\KwOut{Result(aclFailures, bpmnFailures, goalFailures, oracleFailures)}

order BPMN activities topologically and append each activity's effect to the initial SOIL script\;
compile the extended script into checkpoints $\Sigma_0,\dots,\Sigma_n$ ($\Gamma_B$), computing $\mu_{\Sigma_i}$ at every checkpoint\;
\ForEach{checkpoint that is initial or follows an effect}{
  check every ACL/USE invariant\; stop at the first violation\;
}
\ForEach{activity $a_i$}{
  evaluate $\mathit{pre}(a_i)$ at $\Sigma_{i-1}$ and $\mathit{post}(a_i)$ at $\Sigma_i$\; record every violation\;
}
evaluate every root i* goal via $\mu_{\Sigma_n}$, for every $(g,t') \in \mathrm{dom}(\mu_{\Sigma_n})$\; record every unfulfilled goal\;
\ForEach{$(e,t') \in \mathrm{dom}(\mathit{oracle}) \cap \mathrm{dom}(\mu_{\Sigma_n})$}{
  compare $\mu_{\Sigma_n}(e,t')$ against $\mathit{oracle}(e,t')$\; record every mismatch\;
}
\Return Result(aclFailures, bpmnFailures, goalFailures, oracleFailures)\;
\end{algorithm}
 from within a section, or copy the
% relevant \begin{algorithm}...\end{algorithm} block directly. Labels are kept
% as in the original (alg::acl2use, alg::checking); re-check for label clashes
% before including both this file and any other file using the same labels.

% ---------------------------------------------------------------------------
% Originally in Part B (sec:acl2use, "Translating ACL into a Class Diagram"),
% introduced by: "Algorithm~\ref{alg::acl2use} summarizes the translation and
% its compatibility rule: for each concrete role pair $(R_1, R_2)$ sharing a
% group scope, no invariant is generated if they can never share an instance
% or if $\mathit{Compat}$ declares them compatible over that scope; otherwise
% \texttt{inv NoConflict\_$R_1$\_$R_2$} is generated on \texttt{Agent}."
% ---------------------------------------------------------------------------
\begin{algorithm}
\small
\ttfamily
\caption{Translating an ACL model into a USE class model}
\label{alg::acl2use}
\KwIn{an ACL model $(E, R, G, \mathit{Assoc}, \mathit{Compat})$}
\KwOut{a USE \texttt{.use} class model}

create a USE class for every entity in $E$, role in $R$, and group in $G$\;
translate role extension into USE generalization\;
\ForEach{relation in $\mathit{Assoc}$}{
  add the corresponding USE association/aggregation/composition, with the same cardinalities\;
}
create class \texttt{Agent}\;
\ForEach{ownership relation between a group $g$ and a role $r$}{
  add association \texttt{Agent\_plays\_$r$} between \texttt{Agent} and $r$\;
  add association \texttt{$r$\_in\_$g$} between $g$ and $r$, with cardinality $\mathit{card}(g,r)$\;
}
\ForEach{ownership relation between a parent group and a child group}{
  add a composition between them\;
}
\ForEach{pair of concrete roles $(R_1,R_2)$ that can share a group scope}{
  \eIf{$(R_1,R_2) \in \mathit{Compat}$ over that scope}{
    generate no invariant\;
  }{
    generate \texttt{inv NoConflict\_$R_1$\_$R_2$} on \texttt{Agent}\;
  }
}
\Return the generated \texttt{.use} class model\;
\end{algorithm}

% ---------------------------------------------------------------------------
% Originally in Part E (sec:checking, "Checking i* and BPMN over State
% Transitions"), introduced by: "Algorithm~\ref{alg::checking} restates this
% as a procedure," following the four failure-list definitions
% (aclFailures, bpmnFailures, goalFailures, oracleFailures).
% ---------------------------------------------------------------------------
\begin{algorithm}
\small
\ttfamily
\caption{Checking i* against BPMN on a checkpoint trace}
\label{alg::checking}
\KwIn{a USE class model generated from ACL, an initial SOIL script compiled from AOL, an i* model, an ISCN oracle, a BPMN model}
\KwOut{Result(aclFailures, bpmnFailures, goalFailures, oracleFailures)}

order BPMN activities topologically and append each activity's effect to the initial SOIL script\;
compile the extended script into checkpoints $\Sigma_0,\dots,\Sigma_n$ ($\Gamma_B$), computing $\mu_{\Sigma_i}$ at every checkpoint\;
\ForEach{checkpoint that is initial or follows an effect}{
  check every ACL/USE invariant\; stop at the first violation\;
}
\ForEach{activity $a_i$}{
  evaluate $\mathit{pre}(a_i)$ at $\Sigma_{i-1}$ and $\mathit{post}(a_i)$ at $\Sigma_i$\; record every violation\;
}
evaluate every root i* goal via $\mu_{\Sigma_n}$, for every $(g,t') \in \mathrm{dom}(\mu_{\Sigma_n})$\; record every unfulfilled goal\;
\ForEach{$(e,t') \in \mathrm{dom}(\mathit{oracle}) \cap \mathrm{dom}(\mu_{\Sigma_n})$}{
  compare $\mu_{\Sigma_n}(e,t')$ against $\mathit{oracle}(e,t')$\; record every mismatch\;
}
\Return Result(aclFailures, bpmnFailures, goalFailures, oracleFailures)\;
\end{algorithm}
 from within a section, or copy the
% relevant \begin{algorithm}...\end{algorithm} block directly. Labels are kept
% as in the original (alg::acl2use, alg::checking); re-check for label clashes
% before including both this file and any other file using the same labels.

% ---------------------------------------------------------------------------
% Originally in Part B (sec:acl2use, "Translating ACL into a Class Diagram"),
% introduced by: "Algorithm~\ref{alg::acl2use} summarizes the translation and
% its compatibility rule: for each concrete role pair $(R_1, R_2)$ sharing a
% group scope, no invariant is generated if they can never share an instance
% or if $\mathit{Compat}$ declares them compatible over that scope; otherwise
% \texttt{inv NoConflict\_$R_1$\_$R_2$} is generated on \texttt{Agent}."
% ---------------------------------------------------------------------------
\begin{algorithm}
\small
\ttfamily
\caption{Translating an ACL model into a USE class model}
\label{alg::acl2use}
\KwIn{an ACL model $(E, R, G, \mathit{Assoc}, \mathit{Compat})$}
\KwOut{a USE \texttt{.use} class model}

create a USE class for every entity in $E$, role in $R$, and group in $G$\;
translate role extension into USE generalization\;
\ForEach{relation in $\mathit{Assoc}$}{
  add the corresponding USE association/aggregation/composition, with the same cardinalities\;
}
create class \texttt{Agent}\;
\ForEach{ownership relation between a group $g$ and a role $r$}{
  add association \texttt{Agent\_plays\_$r$} between \texttt{Agent} and $r$\;
  add association \texttt{$r$\_in\_$g$} between $g$ and $r$, with cardinality $\mathit{card}(g,r)$\;
}
\ForEach{ownership relation between a parent group and a child group}{
  add a composition between them\;
}
\ForEach{pair of concrete roles $(R_1,R_2)$ that can share a group scope}{
  \eIf{$(R_1,R_2) \in \mathit{Compat}$ over that scope}{
    generate no invariant\;
  }{
    generate \texttt{inv NoConflict\_$R_1$\_$R_2$} on \texttt{Agent}\;
  }
}
\Return the generated \texttt{.use} class model\;
\end{algorithm}

% ---------------------------------------------------------------------------
% Originally in Part E (sec:checking, "Checking i* and BPMN over State
% Transitions"), introduced by: "Algorithm~\ref{alg::checking} restates this
% as a procedure," following the four failure-list definitions
% (aclFailures, bpmnFailures, goalFailures, oracleFailures).
% ---------------------------------------------------------------------------
\begin{algorithm}
\small
\ttfamily
\caption{Checking i* against BPMN on a checkpoint trace}
\label{alg::checking}
\KwIn{a USE class model generated from ACL, an initial SOIL script compiled from AOL, an i* model, an ISCN oracle, a BPMN model}
\KwOut{Result(aclFailures, bpmnFailures, goalFailures, oracleFailures)}

order BPMN activities topologically and append each activity's effect to the initial SOIL script\;
compile the extended script into checkpoints $\Sigma_0,\dots,\Sigma_n$ ($\Gamma_B$), computing $\mu_{\Sigma_i}$ at every checkpoint\;
\ForEach{checkpoint that is initial or follows an effect}{
  check every ACL/USE invariant\; stop at the first violation\;
}
\ForEach{activity $a_i$}{
  evaluate $\mathit{pre}(a_i)$ at $\Sigma_{i-1}$ and $\mathit{post}(a_i)$ at $\Sigma_i$\; record every violation\;
}
evaluate every root i* goal via $\mu_{\Sigma_n}$, for every $(g,t') \in \mathrm{dom}(\mu_{\Sigma_n})$\; record every unfulfilled goal\;
\ForEach{$(e,t') \in \mathrm{dom}(\mathit{oracle}) \cap \mathrm{dom}(\mu_{\Sigma_n})$}{
  compare $\mu_{\Sigma_n}(e,t')$ against $\mathit{oracle}(e,t')$\; record every mismatch\;
}
\Return Result(aclFailures, bpmnFailures, goalFailures, oracleFailures)\;
\end{algorithm}
 from within a section, or copy the
% relevant \begin{algorithm}...\end{algorithm} block directly. Labels are kept
% as in the original (alg::acl2use, alg::checking); re-check for label clashes
% before including both this file and any other file using the same labels.

% ---------------------------------------------------------------------------
% Originally in Part B (sec:acl2use, "Translating ACL into a Class Diagram"),
% introduced by: "Algorithm~\ref{alg::acl2use} summarizes the translation and
% its compatibility rule: for each concrete role pair $(R_1, R_2)$ sharing a
% group scope, no invariant is generated if they can never share an instance
% or if $\mathit{Compat}$ declares them compatible over that scope; otherwise
% \texttt{inv NoConflict\_$R_1$\_$R_2$} is generated on \texttt{Agent}."
% ---------------------------------------------------------------------------
\begin{algorithm}
\small
\ttfamily
\caption{Translating an ACL model into a USE class model}
\label{alg::acl2use}
\KwIn{an ACL model $(E, R, G, \mathit{Assoc}, \mathit{Compat})$}
\KwOut{a USE \texttt{.use} class model}

create a USE class for every entity in $E$, role in $R$, and group in $G$\;
translate role extension into USE generalization\;
\ForEach{relation in $\mathit{Assoc}$}{
  add the corresponding USE association/aggregation/composition, with the same cardinalities\;
}
create class \texttt{Agent}\;
\ForEach{ownership relation between a group $g$ and a role $r$}{
  add association \texttt{Agent\_plays\_$r$} between \texttt{Agent} and $r$\;
  add association \texttt{$r$\_in\_$g$} between $g$ and $r$, with cardinality $\mathit{card}(g,r)$\;
}
\ForEach{ownership relation between a parent group and a child group}{
  add a composition between them\;
}
\ForEach{pair of concrete roles $(R_1,R_2)$ that can share a group scope}{
  \eIf{$(R_1,R_2) \in \mathit{Compat}$ over that scope}{
    generate no invariant\;
  }{
    generate \texttt{inv NoConflict\_$R_1$\_$R_2$} on \texttt{Agent}\;
  }
}
\Return the generated \texttt{.use} class model\;
\end{algorithm}

% ---------------------------------------------------------------------------
% Originally in Part E (sec:checking, "Checking i* and BPMN over State
% Transitions"), introduced by: "Algorithm~\ref{alg::checking} restates this
% as a procedure," following the four failure-list definitions
% (aclFailures, bpmnFailures, goalFailures, oracleFailures).
% ---------------------------------------------------------------------------
\begin{algorithm}
\small
\ttfamily
\caption{Checking i* against BPMN on a checkpoint trace}
\label{alg::checking}
\KwIn{a USE class model generated from ACL, an initial SOIL script compiled from AOL, an i* model, an ISCN oracle, a BPMN model}
\KwOut{Result(aclFailures, bpmnFailures, goalFailures, oracleFailures)}

order BPMN activities topologically and append each activity's effect to the initial SOIL script\;
compile the extended script into checkpoints $\Sigma_0,\dots,\Sigma_n$ ($\Gamma_B$), computing $\mu_{\Sigma_i}$ at every checkpoint\;
\ForEach{checkpoint that is initial or follows an effect}{
  check every ACL/USE invariant\; stop at the first violation\;
}
\ForEach{activity $a_i$}{
  evaluate $\mathit{pre}(a_i)$ at $\Sigma_{i-1}$ and $\mathit{post}(a_i)$ at $\Sigma_i$\; record every violation\;
}
evaluate every root i* goal via $\mu_{\Sigma_n}$, for every $(g,t') \in \mathrm{dom}(\mu_{\Sigma_n})$\; record every unfulfilled goal\;
\ForEach{$(e,t') \in \mathrm{dom}(\mathit{oracle}) \cap \mathrm{dom}(\mu_{\Sigma_n})$}{
  compare $\mu_{\Sigma_n}(e,t')$ against $\mathit{oracle}(e,t')$\; record every mismatch\;
}
\Return Result(aclFailures, bpmnFailures, goalFailures, oracleFailures)\;
\end{algorithm}
